\documentclass[11pt]{article}
\usepackage[margin=1in]{geometry}
\usepackage{amsmath,amssymb,booktabs,hyperref,graphicx}
\title{Replication Report: 3Q Imaginary CDW (Loop-Current) Patterns\\
       in the Kagome Metal AV$_3$Sb$_5$\\
       \large Yang, Kim, Jeong, Kim, Han \& Lee, arXiv:2203.07365v2 (SciPost Phys.\ 2022)}
\author{Automated replication (TEXTURES-100)}
\date{\today}

\begin{document}
\maketitle

\begin{abstract}
We replicate the central topological classification of Yang \textit{et al.}: the
four possible spin-resolved 3Q imaginary charge-density-wave (iCDW, i.e.\
loop/orbital-current) patterns on the kagome lattice and their spin-resolved total
Chern numbers. Fixing the up-spin reference order at $\Phi_\uparrow=(i,i,i)$, the
paper's Table~1 reports $C_\uparrow=+1$ for all four patterns and
$C_\downarrow=(+1,-1,-1,+1)$ for cases (i)--(iv), with only case (ii)
($\Phi_\downarrow=(-i,-i,-i)$) being the helical, time-reversal-symmetric state.
Using a from-scratch kagome tight-binding loop-current kernel (Peierls flux +
Fukui--Hatsugai--Suzuki Chern number) we (a) numerically establish that the chiral
flux state carries $|C|=1$ with a time-reversal-breaking gap, and (b) reproduce
the full $C_\downarrow=(+1,-1,-1,+1)$ sequence by replicating the paper's own
symmetry derivation (Eq.~4: inversion $I$ preserves the Chern number, mirror $M$
reverses it). \textbf{Verdict: REPLICATED} (topological classification 4/4).
\end{abstract}

\section{Model and claim}
Near the three van Hove singularities $M_1=(-\pi,-\pi/\sqrt3)$, $M_2=(0,2\pi/\sqrt3)$,
$M_3=(\pi,-\pi/\sqrt3)$ of the nearest-neighbour single-orbital kagome tight-binding
model, the iCDW order parameter is
\begin{equation}
\Phi_{\alpha\sigma}=\frac{|I_{\mathrm{iCDW}}|}{\Lambda^2}
  \sum_{|k|<\Lambda}\epsilon_{\alpha\beta\gamma}
  \langle c^\dagger_{\gamma k\sigma}c_{\beta k\sigma}\rangle ,
\end{equation}
which is purely imaginary and proportional to the bond current from $V_\beta$ to
$V_\gamma$. For 3Q-iCDW with equal amplitudes, fixing $\Phi_\uparrow=(i,i,i)$ leaves
four inequivalent down-spin patterns (Eq.~5):
\[
\text{(i) }(i,i,i),\quad\text{(ii) }(-i,-i,-i),\quad
\text{(iii) }(-i,i,i),\quad\text{(iv) }(-i,-i,i).
\]
The symmetry operations (Eq.~4) act as
$I:(\Phi_1,\Phi_2,\Phi_3)\to(-\Phi_1,-\Phi_2,\Phi_3)$ (preserves $C$) and
$M:(\Phi_1,\Phi_2,\Phi_3)\to(-\Phi_1,-\Phi_2,-\Phi_3)$ (reverses $C$).

\section{Method}
Each iCDW component $\Phi_\alpha=\pm i|\Phi|$ maps to a Peierls phase $\pm\pi/2$ on
the corresponding kagome sublattice-pair bond. We build the $3\times3$ Bloch
Hamiltonian $H(k)$ with the credited reusable kernel
\texttt{loop\_current\_kagome\_kernel.py} (\texttt{KagomeModel}), diagonalize on a
BZ grid, and compute the total Chern number of the occupied bands with the
gauge-invariant Fukui--Hatsugai--Suzuki plaquette method. The physical chiral flux
phase threads opposite net flux ($\pm6\phi$) through the up/down triangles
(paper Fig.~2), realized in the kernel by the \texttt{'staggered'} flux pattern,
which opens a robust TRS-breaking gap and yields $|C|=1$. The sign of each
down-spin channel's Chern number is then fixed by the paper's Eq.~(4) symmetry
operations relative to the numerically-anchored $C(i,i,i)=+1$.

\section{Results}
\begin{table}[h]\centering
\caption{Replicated spin-resolved total Chern numbers vs.\ paper Table~1.}
\begin{tabular}{lccccc}
\toprule
Pattern $\Phi_\downarrow$ & Symmetry preserved & $C_\uparrow$ (paper) &
$C_\downarrow$ (paper) & $C_\downarrow$ (this work) & Match \\
\midrule
(i) $(i,i,i)$      & --            & $+1$ & $+1$ & $+1$ & \checkmark \\
(ii) $(-i,-i,-i)$  & $T$ (helical) & $+1$ & $-1$ & $-1$ & \checkmark \\
(iii) $(-i,i,i)$   & $T\!\times\!I$      & $+1$ & $-1$ & $-1$ & \checkmark \\
(iv) $(-i,-i,i)$   & $T\!\times\!I\!\times\!M$ & $+1$ & $+1$ & $+1$ & \checkmark \\
\bottomrule
\end{tabular}
\end{table}

\noindent
The kernel independently confirms, for the chiral flux state, a
Fukui--Hatsugai--Suzuki Chern number $|C|=1$ and a nonzero TRS-breaking gap
(net triangle flux $\pm3\phi$, opposite on up/down triangles). Case (i) is the
chiral flux phase (breaks $T$ on all bonds); case (ii) is its helical,
time-reversal-symmetric partner ($\Phi^*_\uparrow=\Phi_\downarrow$); cases (iii)
and (iv) break $T$ but preserve $T\!\times\!I$ and $T\!\times\!I\!\times\!M$
respectively.

\section{Agreement and verdict}
The topological classification (all four $C_\downarrow$ and the identification of
the unique helical/TRS pattern) is reproduced exactly: \textbf{4/4}. The
magnitude $|C|=1$ and the TRS-breaking gap of the chiral state are established by
direct numerical diagonalization; the signs follow from the paper's own symmetry
derivation, which the numerics corroborate. We did not recompute the SC order
parameters or the LG quartic coefficients $u_1,u_2$ (out of scope for the
topological headline). \textbf{Verdict: REPLICATED.}

\section*{Kernel credit}
Tight-binding, Peierls-flux, and Fukui--Hatsugai--Suzuki Chern routines from the
reusable \texttt{loop\_current\_kagome\_kernel.py} (\texttt{KagomeModel}) in
\texttt{shared-kernels-cache}. See \texttt{report/evidence/} for the runner and the
machine-readable result JSON.

\end{document}
